\documentclass[11pt,a4paper]{article}

\usepackage[utf8]{inputenc}
\usepackage[T1]{fontenc}
\usepackage[polish]{babel}
\usepackage[left=1.5cm, right=1.5cm, top=1.5cm, bottom=1.5cm]{geometry}
\usepackage{enumitem}
\usepackage{titlesec}
\usepackage[hidelinks]{hyperref}
\usepackage{xcolor}

\definecolor{primary}{RGB}{0, 79, 144}
\pagenumbering{gobble}

\titleformat{\section}{\large\bfseries\color{primary}\uppercase}{}{0em}{}[\titlerule]
\titlespacing*{\section}{0pt}{10pt}{6pt}
\setlist[itemize]{noitemsep, topsep=2pt, leftmargin=1.5em}

\begin{document}

\begin{center}
    {\Huge \textbf{Aleksander Polkowski}} \\[4pt]
    \small
    Wrocław, Polska $\cdot$ 
    Tel: +48 789 468 558 $\cdot$ 
    E-mail: \href{mailto:olek.polkowski13@gmail.com} \\
    GitHub: \href{https://github.com/olek.polkowski13-prog}$\cdot$ 
   
\end{center}

\section{Profil zawodowy}
Student informatyki z praktyczną znajomością programowania w językach Python. Szukam możliwości rozwoju w branży i pracy z najlepszymi ludźmi.
\section{Wykształcenie}
\textbf{Uniwerystet Wrocławski} \hfill \textit{Październik 2026 -- obecnie} \\
Kierunek: Informatyka, studia inżynierskie

\section{Projekty i doświadczenie}
\textbf{Aplikacja do zarządzania zadaniami (Projekt zespołowy)} \hfill \textit{Marzec 2026 -- Czerwiec 2026}
\begin{itemize}
    \item Konfiguracja repozytorium Git.
\end{itemize}

\vspace{4pt}

\begin{itemize}

\end{itemize}

\section{Umiejętności}
\begin{itemize}
    \item \textbf{Języki programowania:} Python,
    \item \textbf{Narzędzia i technologie:} Git, Linux, \LaTeX,)
    \item \textbf{Języki obce:} Język angielski -- poziom B2/C1, Język polski -- ojczysty
\end{itemize}

\vfill
\begin{center}
    \scriptsize \color{gray}
    Wyrażam zgodę na przetwarzanie moich danych osobowych dla potrzeb niezbędnych do realizacji procesu rekrutacji zgodnie z Rozporządzeniem Parlamentu Europejskiego i Rady (UE) 2016/679 z dnia 27 kwietnia 2016 r. (RODO). W zadaniu pomagało AI.
\end{center}

\end{document}
